\chapter{Sprint 2 — Video Indexing and Keyword Search}

\section*{Introduction}
\addcontentsline{toc}{section}{Introduction}

Sprint 2 builds the core of the Videos Search API: ingestion of video documents into OpenSearch
with normalized metadata, keyword (BM25) search over weighted fields, and aggregation-based
facets that power the portal's filter sidebar.

%% ─── Sprint Backlog ─────────────────────────────────────────────────────────
\section{Sprint Backlog}

\begin{table}[H]
\centering
\caption{Sprint 2 Backlog}
\label{tab:sprint2_backlog}
\begin{tabularx}{\textwidth}{clXcc}
\toprule
\textbf{ID} & \textbf{Story} & \textbf{Task} & \textbf{Effort (d)} & \textbf{Prio.} \\
\midrule
T2.1 & US-08 & Flask + Flask-RESTX scaffolding with OpenSearch client injection & 1 & M \\
T2.2 & US-08 & Video document schema (title, description, tags, metadata, orientation) & 1 & M \\
T2.3 & US-08 & Index mapping + read/write aliases (\texttt{videos\_read}, \texttt{videos\_write}) configuration & 1 & M \\
T2.4 & US-08 & Ingestion endpoint with Marshmallow validation                        & 1 & M \\
T2.5 & US-08 & Keyword search (\texttt{multi\_match}) with per-field boosts           & 1 & M \\
T2.6 & US-08 & Filter clauses (categories, tags, duration, orientation, resolution)  & 1 & M \\
T2.7 & US-08 & Facet aggregation endpoint for the filter sidebar                     & 0.5 & M \\
T2.8 & US-08 & Latest-videos and categories listing endpoints                        & 0.5 & S \\
T2.9 & US-15 & Dockerization of the service                                          & 0.5 & M \\
\bottomrule
\end{tabularx}
\end{table}

%% ─── Design ─────────────────────────────────────────────────────────────────
\section{Design}

\subsection{Use Case Diagram for Sprint 2}

The video-search portal's anonymous visitor drives this sprint: browse, search, filter, and view
video detail. These cases belong to the Video Search Portal boundary of the global use-case
diagram (Figure~\ref{fig:uc_global} in Chapter~2).

\subsection{Class Diagram for Sprint 2}

The OpenSearch document model used by Sprint 2 maps each video to one document with the fields
listed in Table~\ref{tab:video_doc}. These fields are denormalized for fast lexical search and
aggregation.

\begin{table}[H]
\centering
\caption{Video document schema in OpenSearch}
\label{tab:video_doc}
\small
\begin{tabularx}{\textwidth}{lX}
\toprule
\textbf{Field} & \textbf{Purpose} \\
\midrule
\texttt{video\_id}             & Unique identifier (used as primary boost target). \\
\texttt{title}                 & Primary textual field for search; boosted $\times 3$. \\
\texttt{description}           & Long text; boosted $\times 2$. \\
\texttt{tags}                  & Multi-valued keywords; boosted $\times 2$. \\
\texttt{categories}            & Multi-valued categories; boosted $\times 2$. \\
\texttt{location}, \texttt{resolution}, \texttt{orientation}, \texttt{owner\_name} & Metadata fields used for filters. \\
\texttt{duration\_sec}         & Duration in seconds; used by duration range filter. \\
\texttt{created\_at}, \texttt{views\_max} & Sort keys for newest/oldest/views sort modes. \\
\texttt{rms.data.*}, \texttt{cts.data.*} & Legacy nested metadata kept as fallback fields. \\
\bottomrule
\end{tabularx}
\normalsize
\end{table}

\subsection{Sequence Diagram: Video Indexing}

Figure~\ref{fig:seq_video_indexing} shows how an ingestion job flows through the Videos Search
API into OpenSearch.

\begin{figure}[H]
\centering
\begin{tikzpicture}[x=1cm, y=1cm, font=\fontsize{7}{8}\selectfont]
\node[sdWork] (pIngest) at ( 0.0, 0) {Ingestion\\Worker};
\node[sdBE]   (pVSAPI)  at ( 3.2, 0) {Videos Search\\API};
\node[sdExt]  (pEmbed)  at ( 6.4, 0) {Embedding\\API};
\node[sdData] (pOS)     at ( 9.6, 0) {OpenSearch};

\foreach \x in {0.0, 3.2, 6.4, 9.6}
  \draw[dashed, gray!50] (\x,-0.3) -- (\x,-8.8);

\draw[sdmsg] (0.0,-0.85) -- node[above, pos=0.5]{new video payload} (3.2,-0.85);
\draw[sdmsg] (3.2,-1.57) -- node[above, pos=0.5]{validate + normalize fields} (3.2,-1.57);
\draw[sdmsg] (3.2,-2.29) -- node[above, pos=0.5]{embed(title + tags + transcript)} (6.4,-2.29);
\draw[sdret] (6.4,-3.01) -- node[above, pos=0.5]{vector [d dimensions]} (3.2,-3.01);
\draw[sdmsg] (3.2,-3.73) -- node[above, pos=0.5]{PUT /videos/\{id\} with metadata + vector} (9.6,-3.73);
\draw[sdret] (9.6,-4.45) -- node[above, pos=0.5]{201 created / updated} (3.2,-4.45);
\draw[sdmsg] (3.2,-5.17) -- node[above, pos=0.5]{refresh/alias update (batch policy)} (9.6,-5.17);
\draw[sdret] (9.6,-5.89) -- node[above, pos=0.5]{index ready for queries} (3.2,-5.89);
\draw[sdret] (3.2,-6.61) -- node[above, pos=0.5]{ingestion status=SUCCESS} (0.0,-6.61);

\node[sdnote] at (11.9,-4.1) {Document stores\\text fields +\\embedding vector};
\draw[->, dashed, gray!50] (9.6,-4.1) -- (11.0,-4.1);
\end{tikzpicture}
\caption{Video indexing sequence (metadata + embedding upsert)}
\label{fig:seq_video_indexing}
\end{figure}

%% ─── Implementation ────────────────────────────────────────────────────────
\section{Implementation}

\subsection{Flask + Flask-RESTX with OpenSearch Client Injection}

The Videos Search API uses \textbf{Flask} with \textbf{Flask-RESTX} resource classes. The
OpenSearch client is injected per-request via Flask's application context (\texttt{g}). Input is
validated with \textbf{Marshmallow} schemas. Flask was kept for compatibility with the existing
\textbf{Python 3.8} runtime and the Elasticsearch-compatible OpenSearch Python client.

Key endpoints exposed by this sprint:

\begin{itemize}
  \item \texttt{POST /api/videos/query} — Multi-mode search (semantic, filters, or auto).
  \item \texttt{POST /api/videos/vsearch} — Full-text vector semantic search.
  \item \texttt{GET /api/videos/categories} — Fetch available categories.
  \item \texttt{GET /api/videos/getlatest} — Latest video clips.
  \item \texttt{GET /api/videos/facets} — Aggregation facets for the filter sidebar.
\end{itemize}

\subsection{Indexing in OpenSearch}

When a video is ingested, the API normalizes metadata (title, category, tags, duration,
orientation, source platform) and writes one document to the OpenSearch index. This gives fast
keyword search (BM25), filters, and facet aggregations for the portal sidebar. The index is
accessed through dedicated aliases: \texttt{videos\_write} for ingestion updates and
\texttt{videos\_read} for query traffic, which keeps reindexing operations safe and
zero-downtime.

\subsection{Keyword Search (BM25)}

Keyword search is the traditional full-text path. It relies on \textbf{exact and fuzzy term
matches} between the user query and indexed text fields. At query time, the API issues a
\texttt{multi\_match} query that scores each document using \textbf{BM25}, weighted per field so
stronger signals dominate.

\begin{table}[H]
\centering
\caption{Field weights used by the keyword ranker}
\begin{tabular}{ll}
\toprule
\textbf{Field} & \textbf{Boost} \\
\midrule
\texttt{video\_id}                                  & $\times 4$ \\
\texttt{title} / \texttt{rms.data.name}             & $\times 3$ \\
\texttt{description} / \texttt{tags} / \texttt{categories} & $\times 2$ \\
\texttt{location} / \texttt{resolution} / \texttt{orientation} / \texttt{owner\_name} & $\times 1$ \\
\bottomrule
\end{tabular}
\end{table}

Keyword search is fast, predictable, and great for precise queries (exact titles or IDs). It fails
when users describe intent rather than vocabulary (``a crowd celebrating a football goal''), which
motivates the semantic retrieval built in Sprint 3.

\begin{figure}[H]
\centering
\begin{tikzpicture}[x=1cm,y=1cm,
  box/.style={draw=teal!60!black, fill=teal!8, rounded corners=4pt,
              minimum width=3.4cm, minimum height=0.95cm, align=center, font=\small},
  data/.style={draw=violet!60!black, fill=violet!8, rounded corners=4pt,
               minimum width=3.6cm, minimum height=0.95cm, align=center, font=\small},
  ext/.style={draw=gray!60, fill=gray!8, rounded corners=4pt,
              minimum width=3.0cm, minimum height=0.95cm, align=center, font=\small},
  arr/.style={->, >=Stealth, line width=0.9pt}
]
\node[ext]  (portal) at (-5.6, 0)   {User query\\(text + filters)};
\node[box]  (api)    at (-1.8, 0)   {Videos Search API\\analyze tokens};
\node[box]  (mm)     at ( 2.2, 0)   {\texttt{multi\_match}\\+ filter clauses};
\node[data] (os)     at ( 6.3, 0)   {OpenSearch\\inverted index (BM25)};
\node[data] (hits)   at ( 2.2,-2.5) {Ranked hits\\\_score per doc};

\draw[arr] (portal) -- (api);
\draw[arr] (api) -- (mm);
\draw[arr] (mm) -- (os);
\draw[arr] (os.south) |- (hits.east);
\draw[arr] (hits.west) -| node[below, font=\footnotesize, pos=0.25]{results} (portal.south);
\end{tikzpicture}
\caption{Keyword (BM25) search flow in the Videos Search API}
\label{fig:keyword_search}
\end{figure}

\subsection{Facets and Filter Sidebar}

The facet aggregation endpoint returns, for the current result set, the distinct values and
counts for \texttt{categories}, \texttt{tags}, \texttt{resolutions}, and \texttt{orientation}.
The portal uses these numbers to render the filter sidebar shown in
Figure~\ref{fig:screen_vp_search}.

\subsection{Delivered Interfaces}


\begin{figure}[H]
\centering
\includegraphics[width=\textwidth]{\detokenize{figures/videosearch/browse page.png}}
\caption{Video Search Portal — search results with filters}
\label{fig:screen_vp_search}
\end{figure}

\begin{figure}[H]
\centering
\includegraphics[width=\textwidth]{\detokenize{figures/videosearch/videodetails page.png}}
\caption{Video Search Portal — video detail page}
\label{fig:screen_vp_detail}
\end{figure}

\section*{Conclusion}
\addcontentsline{toc}{section}{Conclusion}

Sprint 2 produced a working video-search stack that serves keyword and facet-driven browsing.
The next sprint adds semantic retrieval, parallel hybrid search, and a position-aware ranker on
top of this lexical baseline.
